\subsection{Steady-State Resource Profiling (Scenarios 1 and 2)}
To evaluate the static computational and memory footprint of the telemetry infrastructure (Hypothesis $H_{1,3}$) without introducing temporal pseudoreplication (\textbf{ADR-0003}), the profiler scripts (\texttt{run-scenario1-baseline.py} and \texttt{run-scenario2-overhead.py}) sampled the cluster state across $N=30$ independent observations at a strict $16\,\text{s}$ interval over a $480\,\text{s}$ steady-state window. Because the Kubernetes \texttt{metrics-server} scrapes kubelet \texttt{cAdvisor} counters every $15\,\text{s}$, enforcing a $16\,\text{s}$ polling cadence guarantees that every recorded sample reflects a fresh, uncached physical measurement \cite{hurlbert_1984, jain_1991}.

Computational values were algorithmically normalized: CPU utilization was converted to uniform millicore units ($\text{mCPU}$), while resident working-set memory allocations were standardized to Mebibytes ($\text{MiB}$). In Scenario~1, the profiler captured the baseline workload demand ($R_{\text{S1}}$) of the scoped OpenTelemetry Astronomy Shop (\texttt{otel-demo}, \textbf{ADR-0005}) running on the Vanilla K3d cluster without telemetry agents. In Scenario~2, the script concurrently aggregated metrics across both the application namespace (\texttt{otel-demo}, capturing $R_{\text{workload}}$ under active OpenTelemetry OTLP export) and the observability namespace (\texttt{monitoring}, isolating $R_{\text{PANOPTES}}$ across Prometheus, Grafana Loki, Grafana Tempo, Grafana, and the OpenTelemetry Collector), persisting the time-series records to \texttt{scenario1\_workload\_baseline.csv} and \texttt{scenario2\_overhead\_panoptes.csv}.
